\chapter{Limitations and Future Research}
\label{ch:limitations}

\section{The Inferential Scope of the Current Design}
\label{sec:inferential-scope}
This thesis was designed to answer a narrow question with precision: what textual and structural patterns characterize cross-state diffusion in the first post-\textit{Dobbs} legislative cycle, and which institutional and coalitional features predict them? The findings reported in Chapter~\ref{ch:findings}, and the theoretical apparatus advanced in Chapter~\ref{ch:discussion}---constitutional vacuum diffusion, reactive template mobilization, and the vertical--horizontal coalition architecture distinction ---carry implications that extend well beyond that narrow question, and present chapter accounts for the distance between the two. The computational design rigorously establishes the patterns it was built to detect; it leaves a corresponding set of mechanism-level questions for subsequent investigation. This chapter identifies that distance candidly, specifies the limits it imposes on the claims advanced in Chapter~\ref{ch:discussion}, and articulates the future research agenda that would close it.

The organizing premise is the one developed at the close of Chapter~\ref{ch:discussion}: different research instruments answer different questions, and the textual and documentary materials assembled for this thesis answer rigorously the question they were designed to answer while bracketing, for subsequent investigation, those they were not. The sections that follow proceed from the narrower methodological limits---what TF--IDF cosine similarity can and cannot answer about legislative coordination (\S\ref{sec:limits-computational})---through the specification-level limits of the dyadic regression (\S\ref{sec:limits-dyadic}) and the scope limits imposed by a single-year, single-domain corpus (\S\ref{sec:limits-scope}), to an explicit accounting of demonstration versus inference (\S\ref{sec:accounting}). The chapter then turns to the future research agenda that these limits collectively define (\S\ref{sec:future-research}), centered on direct qualitative inquiry into the drafting processes internal to legislative offices and advocacy organizations. A brief concluding section (\S\ref{sec:designed-scope}) reprises the framing under which these limits are a consequence of methodological prioritization rather than evidentiary incompleteness.

\section{Limitations of the Computational Approach}
\label{sec:limits-computational}

\subsection{Textual Similarity as Necessary but Not Sufficient Evidence of Coordination}
\label{subsec:text-necessary}
The most consequential limit of the computational design concerns the inferential reach of textual similarity itself. TF--IDF vectorization over the 104{,}445 unique n-grams in the 2023 corpus, combined with cosine distance computed over the resulting 188$\times$188 similarity matrix, generates a rigorous numerical description of how closely any two bills converge in surface linguistic content. What that description cannot unambiguously establish is the causal pathways through which the convergence arose. High cosine similarity between two bills is consistent with multiple generative processes, not all of which correspond to the organizational coordination that Chapter~\ref{ch:discussion} theorizes.

Three rival mechanisms are not eliminated by the current design and deserve explicit enumeration. The first is parallel independent drafting by coalitionally aligned legislators working from common doctrinal sources. Legal memoranda circulated through academic networks, law-review commentary on post-\emph{Dobbs} statutory design, and the public briefs filed in ongoing litigation constitute a shared informational environment from which legislators in different states could draw nearly identical statutory language without any direct organizational contact. The textual signature of such parallel drafting would be indistinguishable from the signature of coordinated template deployment, yet the two cases imply different theories of how diffusion operates under a constitutional vacuum.

The second is informal staff copying from academic and advocacy publications. Legislative staff in 2023 had access to an unprecedented volume of publicly available analysis---white papers from think tanks, policy briefs from advocacy organizations, model statutory language published in law reviews, and commentary from legal scholars---any of which could have served as a drafting source without triggering the organizational coordination mechanism the framework identifies. Here too, the textual convergence is genuine, but its causal locus sits outside the drafting infrastructure of the anchor organizations.

The third is post-enactment convergence produced by drafting conventions rather than intentional coordination. Revisor's offices, legislative counsel, and uniform-code influences impose stylistic regularities on statutory language that can generate similarity between bills whose authors never consulted one another or any common organizational source. The residual effect is small in most domains, but on contested, heavily lawyered policy questions, it can be nontrivial. The current design cannot isolate this drafting-conventions component of the observed similarity.

None of these alternatives is eliminated by the evidence assembled for this thesis, and none can be eliminated by any purely textual analysis, however methodologically refined. Discriminating among them requires evidence about the drafting process itself: which sources staff actually consulted, whether and how affiliated organizations circulated language, and at what point in
the drafting cycle particular phrasings entered the statutory text. That
evidence is not available in the corpus.

The implication for the Chapter~\ref{ch:discussion} argument is not that the reactive template mobilization mechanism is wrong, but that the computational design establishes it only by pattern-matching: the signature the framework predicts---compressed simultaneous adoption, horizontal structure, coalition asymmetry, organizational co-presence as the strongest predictor at the dyadic level---is the signature the data exhibit, but a signature is not a mechanism. The framework's mechanism-level claims are therefore best read as one plausible account consistent with the textual and documentary evidence, not as the only account that evidence can support. Comparative weighting among the three rival mechanisms enumerated above and the reactive template mobilization account would require process-level evidence---which sources staff actually consulted, through which channels coordinated language was circulated---that the computational and documentary materials assembled here cannot supply. That comparative weighting is therefore deferred to the qualitative extension specified in Subsection~\ref{subsec:future-qualitative} rather than asserted on the basis of pattern evidence alone.

\subsection{Novelty as Textual Signature, Not Normative Measure of Policy Innovation}
\label{subsec:novelty-signature}

Chapter~\ref{ch:discussion} flagged, in the context of the professionalism reinterpretation, that text-based novelty measures should not be read as measures of policy innovation in the normative sense. The point warrants development here as a methodological caveat that applies across the entire analysis rather than only to the professionalism finding.

The cross-state novelty score constructed for this thesis measures a specific quantity: the textual distance between a bill and its nearest counterpart in another state's 2023 corpus. A low score indicates proximity to existing legislative language elsewhere; a high score indicates textual distance from it. Neither direction carries an implicit evaluation of legislative quality, policy efficacy, or substantive originality. A state that adopts an externally drafted bill well-calibrated to its local institutional and constitutional context may be engaging in more sophisticated policymaking than a state that drafts comparable legislation \emph{de novo} without the benefit of prior statutory experimentation elsewhere. The novelty score is silent on that distinction.

This caveat bears directly on the interpretation of the state-level patterns documented in Chapter~\ref{ch:findings}. Higher novelty scores indicate greater textual distance from out-of-state counterparts; lower scores indicate proximity. That directional convention is purely descriptive: it carries no evaluative content. A state whose bills register low novelty---North Carolina or Tennessee, for example---may be drafting policy of higher quality than a state whose bills register high novelty, because the templates available for adaptation in mature advocacy domains have typically been refined through legal review and litigation testing across multiple jurisdictions. Conversely, high novelty in states with limited infrastructure access may reflect informational isolation rather than creative drafting capacity. Whether any bill in the corpus represents good, durable, or efficacious policy is a question the text-based measure is not designed to answer.

The broader point is that the convergence between substantive policy quality and textual novelty is genuinely weak in domains with dense model-bill infrastructure. High-quality policy often travels through templated channels because the templates have been refined through legal review, litigation testing, and iterative drafting across multiple jurisdictions. Low novelty scores in such domains may therefore track policy sophistication as much as or more than they track unreflective copying. Conversely, high novelty scores in states with limited infrastructure access may reflect informational isolation rather than creative drafting capacity. The interpretive upshot is a distinction between measurement and concept. As \emph{measurement}, the novelty score is well-defined: it captures textual distance between a bill and its nearest cross-state counterpart, with high internal consistency and full reproducibility. As a proxy for the broader \emph{concept} of policy innovation---the substantive originality, doctrinal soundness, implementation viability, judicial resilience, or distributional calibration of legislative output---the score is at best partial and at worst misleading. This thesis treats the textual measure as evidence of diffusion structure, not as a normative ranking of legislative output. Evaluative claims about policy quality would require criteria that lie entirely outside the computational design.

\subsection{Representation Limits of TF--IDF and the N-gram Vector Space}
\label{subsec:tfidf-limits}
The TF--IDF representation adopted in Chapter~\ref{ch:methodology} offers substantial advantages for diffusion analysis: it is transparent, reproducible, interpretable at the feature level, and consistent with the broader text-as-data literature on legislative similarity \parencite{wilkersonTracingFlowPolicy2015, linderTextPolicyMeasuring2020}. It also imposes specific representational costs that any assessment of the thesis's findings should register.

The n-gram vector space treats legislative text as a bag of terms weighted by frequency and inverse document frequency. It discards the semantic equivalence between differently phrased clauses that accomplish the same doctrinal work. Two provisions that protect a clinic from out-of-state subpoenas using different terminology will register as textually dissimilar under TF--IDF even when they are functionally identical. Conversely, two provisions that use identical boilerplate language but produce different substantive effects (because they sit within different enforcement architectures) will register as highly similar. The measure tracks surface linguistic proximity, and in most cases---the abortion corpus included---surface proximity correlates strongly with substantive proximity, but the correspondence is imperfect.

Transformer-based representations such as Legal-BERT and its successors offer partial redress to this limitation by encoding text into embedding spaces that capture semantic equivalence across different surface wordings \parencite{chalkidis-etal-2020-legal}. Replication of the Chapter~\ref{ch:findings} results using semantic embeddings would test whether the coalition asymmetry, the Vermont hub pattern, and the Nevada--Washington near-replication survive the transition to a more semantically sensitive representation. The expectation, given the coarse linguistic convergence the TF--IDF analysis already detects, is that the substantive findings would remain stable and that some additional convergence, invisible to n-gram methods, would emerge. Whether that expectation holds is an open empirical question that the current design does not answer. \S\ref{subsec:future-methodological} returns to semantic replication as a priority methodological extension.

\section{Limitations of the Dyadic Regression Specification}
\label{sec:limits-dyadic}

\subsection{Organizational Presence Versus Drafting Collaboration}
\label{subsec:jaccard-limits}
The dyadic regression reported in Chapter~\ref{ch:findings} established that the Jaccard index of shared PAC donor organizations between two states is a robust and quantitatively comparable predictor of cross-state textual similarity ($\beta = +0.039$, $p < 0.001$, $N = 1{,}035$). The interpretation advanced in Chapter~\ref{ch:discussion} read this result as evidence that organizational co-presence rather than financial transfer drives the observed convergence. That interpretation is the most parsimonious reading of the coefficient, but the measure it relies on warrants explicit qualification.

The Jaccard index registers whether the same organizations maintain an advocacy footprint in both states of a given dyad. It does not register the intensity of that footprint, the specific activities conducted through it, or the presence or absence of actual drafting collaboration between the organization's affiliates in the two states. Two states sharing Planned Parenthood as a DIME-recorded organizational presence might differ substantially in whether that presence involves active legislative drafting assistance, passive lobbying, electoral support for down-ballot candidates, or clinical service provision unconnected to legislative activity. The Jaccard measure does not distinguish among these possibilities, and the dyadic regression cannot therefore identify which forms of organizational co-presence generate the convergence signature.

This limit affects the mechanism-level interpretation more directly than it affects the pattern-level finding. The pattern finding---that dyads sharing organizational footprints exhibit higher textual similarity---is robust to the interpretive question and holds regardless of which activities within those footprints do the causal work. The mechanism finding---that the responsible activity is template circulation through affiliate networks rather than electoral mobilization or clinical advocacy---is the more ambitious claim, and it rests on the theoretical plausibility of the informational-subsidy account in combination with the pattern evidence, not on direct measurement of drafting collaboration. Richer operationalizations
 of advocacy infrastructure---lobbying disclosure filings, IRS Form 990 data on organizational staffing and program expenditures, state-level registration records for legislative liaisons---would provide a more granular test of which components of organizational presence predict convergence. The current design incorporates none of these, and their incorporation is identified in \S\ref{subsec:future-archival} as a priority extension.

 \subsection{Temporal Compression and the Absence of Sequence Data}
 \label{subsec:temporal-compression}
The 2023 corpus captures legislative enactments within a single calendar year. This compression is analytically deliberate---Chapter~\ref{ch:methodology} specified the scope as the first full legislative year after the constitutional vacuum opened---but it imposes a specific inferential cost: the dyadic analysis cannot observe enactment sequencing within the cycle. Whether Vermont SB~37 textually preceded Colorado SB~188 in the sense of being the template from which Colorado drafted, or whether both bills emerged from a common affiliate drafting process without any directional flow between them, or whether Colorado's earlier introduction but later enactment reverses the chronology implied by final passage dates, cannot be determined from
the cross-sectional data alone.

The substantive consequence is that the exporter--importer--reinforcer typology revisited in \S\ref{sec:exporter-importer-reinforcer} cannot, on the present evidence, support directional inference about which states function as origins and which as adopters. What the cross-sectional similarity matrix supports is a structural classification: states whose bills sit at high textual centrality within the network can be identified as candidate hubs, and states whose bills sit at high cross-state novelty can be identified as candidate origins or candidate isolates, but the data cannot adjudicate between these possibilities. The typology is therefore best understood as a pattern-level diagnostic that subsequent enactment-timing analysis would either validate as directional or revise as descriptive. The Vermont hub question identified in \S\ref{sec:exporter-importer-reinforcer} is, at present, underdetermined: it is consistent with Vermont functioning as the originating node in a hub-and-spoke protective network, and it is equally consistent with Vermont, California, New York, and Washington operating as co-drafters within a distributed process whose textual center happens to have been logged as Vermont SB~37 for reasons unrelated to origination. Enactment-timing models on a longitudinal corpus, discussed in \S\ref{subsec:future-longitudinal}, would supply the directional leverage the current design lacks.

\subsection{Statistical Power for Interaction Tests}
\label{subsec:power-limits}
The dyadic dataset comprised 1{,}035 state pairs, which is adequate for robust point estimation of the main effects reported in Table~\ref{tab:dyadic_regression} but limits the statistical power available for interaction tests. Several theoretically interesting questions raised by Chapter~\ref{ch:discussion}---whether the Both Protective effect is conditional on the density of the advocacy infrastructure in the dyad, whether the coalition asymmetry attenuates at higher PAC Jaccard values, and whether professionalism moderates the coalition effects---would require interaction specifications whose power on 1{,}035 dyads is marginal.

Interactions were not included in the preferred specification precisely because the estimated coefficients on plausible interaction terms were unstable across model perturbations, and reporting them as findings would have overstated what a single-year corpus can support.

A longitudinal extension to 2024 and 2025 enactment cohorts, even under conservative assumptions about the number of eligible bills in each subsequent year, would roughly triple the dyadic sample and substantially improve interaction-test power. Whether the coalition asymmetry documented in  this thesis is an unconditional feature of the post-\emph{Dobbs} legislative environment or whether it is conditional on specific organizational and institutional configurations is a question the current corpus lacks the statistical resolution to answer with confidence.

\subsection{Binary Measurement of Advocacy Infrastructure}
\label{subsec:binary-measurement}
The Jaccard construction treats organizational presence within each state as a binary set membership: an organization is either in a state's DIME-recorded advocacy ecosystem or it is not, with no weighting for intensity of engagement. This binary treatment discards information about how deeply each organization is embedded in a given state's political infrastructure. A state where Planned Parenthood maintains five affiliated clinics, a state legislative liaison, a registered lobbyist, and an active candidate-endorsement operation contributes the same Jaccard-index term as a state where Planned Parenthood's only DIME-recorded activity is a single contribution to a gubernatorial campaign. The substantive difference between these two cases is considerable, and the measure does not capture it.

A richer operationalization would weight organizational presence by staffing, lobby-disclosure activity, IRS reporting on program expenditures, or the number of affiliated entities registered in each state. Such weightings would require additional data-collection work---particularly on lobby disclosure records, which are maintained at the state level under inconsistent reporting standards---but would produce a more discriminating test of whether the convergence signature is driven by deep organizational embedding or by mere statutory presence. The current binary measure is adequate to establish the main effect but not to characterize the distributional properties of organizational influence on textual convergence.

A second methodological limitation concerns the clustering specification. The reported standard errors derive from one-way clustering on \texttt{state\_i},
the first state listed in each dyad. For dyadic data with potential within-state correlation operating on both members of each pair, two-way clustering on \texttt{state\_i} and \texttt{state\_j} simultaneously
\parencite{robustinference,aronow2015clusterrobustvarianceestimationdyadic} provides a more rigorous treatment of the dyadic dependence structure. Re-estimation under two-way clustering is a priority for replication. The substantive
results reported in Section 4.9 are unlikely to change in direction or significance under that specification given the conservatism of the Bell--McCaffrey CR2 correction already applied, but the inferential statement will be technically more defensible. 

\section{Limitations of Scope and Selection}
\label{sec:limits-scope}

\subsection{Single Legislative Year}
\label{subsec:single-year}

The corpus captures enacted legislation from a single calendar year---2023, the first full legislative year after \emph{Dobbs}. This temporal window was chosen deliberately: it isolates the initial legislative response to the constitutional vacuum before the feedback effects of early enactments, judicial challenges, and electoral consequences could contaminate the diffusion signal. The cost of that isolation is that the findings describe a moment-specific equilibrium rather than a durable structural feature of post-\emph{Dobbs} policymaking.

Chapter~\ref{ch:discussion} advanced several predictions that are testable only through longitudinal extension. The constitutional vacuum diffusion framework predicts that the coalition asymmetry should attenuate over successive cycles as protective coalitions move beyond initial codification and restrictive coalitions develop new litigation-tested templates for the expanded policy frontier. The reactive template mobilization mechanism predicts that the compressed, simultaneous adoption structure should give way to more sequential, learning-based patterns as the post-shock window matures. Neither prediction can be assessed against a single-year corpus, and both remain stipulated implications of the framework rather than empirically confirmed dynamics. Whether the 2023 patterns represent the leading edge of a transitional regime or a durable reconfiguration of abortion-policy diffusion is a question that only 2024, 2025, and subsequent legislative cycles can answer.

The single-year window also forecloses analysis of pre-\emph{Dobbs} baseline patterns. The constitutional vacuum diffusion framework claims that the 2023 record departs from what established diffusion models would predict; yet the thesis does not construct the pre-\emph{Dobbs} baseline against which that departure could be measured directly. Trigger laws enacted in 2022, the incremental restrictions adopted under \emph{Casey}'s undue-burden framework, and the pre-\emph{Dobbs} protective codifications in states like New York and Illinois constitute an available comparison corpus that the current design does not exploit. Future work incorporating pre-\emph{Dobbs} legislative text would provide a more rigorous before-and-after test of whether the diffusion regime the framework describes is genuinely novel or whether it intensifies patterns already present in the pre-shock period.

\subsection{Single Policy Domain}
\label{subsec:single-domain}

The thesis analyzes abortion policy exclusively. The scope conditions articulated in \S\ref{sec:scope-conditions} specify that the constitutional vacuum diffusion framework should apply to policy domains characterized by predictable preemption collapse, high salience, dense pre-existing advocacy infrastructure, and bilateral mobilization. Those conditions are met in the abortion domain; whether they are met in candidate comparison domains---firearms regulation after \emph{Bruen}, voting rights after \emph{Shelby County}, a hypothetical reversal of \emph{Obergefell}---is argued by analogy rather than established by evidence.

Cross-domain replication would test the framework's most ambitious claim: that constitutional vacuum diffusion is a generalizable regime rather than a description of one policy area's idiosyncratic post-shock dynamics. Without such replication, the possibility remains that the patterns documented here are products of features specific to abortion politics---the depth and duration of the pre-\emph{Roe} organizational infrastructure, the intensity of the morality-policy valence, the uniquely bilateral character of the \emph{Dobbs} opening---rather than products of the structural conditions the framework identifies. The framework's theoretical value depends on its capacity to travel beyond the case that generated it, and establishing that capacity requires empirical work the current design does not undertake.

\subsection{Enacted Bills Only}
\label{subsec:enacted-only}

The corpus includes 188 enacted bills across 46 states. It excludes bills that were introduced but failed to pass, bills that passed one chamber but stalled in the other, and bills that were withdrawn, substituted, or merged into omnibus vehicles before final passage. This selection criterion was methodologically motivated: enacted legislation represents the completed output of the diffusion process and avoids the interpretive ambiguity of bills whose failure may reflect strategic withdrawal, procedural timing, or partisan gatekeeping rather than substantive rejection.

The cost is that the analysis observes only the survivors of the legislative process, which introduces a form of selection bias whose direction is predictable but whose magnitude is unknown. Bills that closely replicated out-of-state templates may have failed precisely because their lack of local adaptation made them politically or legally vulnerable---in which case the enacted corpus underestimates the degree of cross-state templating that was \emph{attempted} in 2023. Alternatively, bills that failed may have been more idiosyncratic than those that passed---in which case the enacted corpus overestimates the baseline novelty of the full introduced population. The direction of the bias depends on whether template fidelity helps or hinders enactment in a given state's legislative environment, and the current design cannot adjudicate.

An extension incorporating the introduced-but-failed population would provide leverage on this question and would also enable analysis of the gatekeeping function: which types of templates survive the legislative process and which are filtered out, and whether that filtering operates differently in restrictive and protective states. LegiScan's comprehensive bill-tracking data make such an extension feasible in principle, though the volume of introduced abortion-related legislation in 2023---substantially larger than the enacted corpus---would require additional data-collection and classification work.

\section{What the Design Can and Cannot Demonstrate}
\label{sec:accounting}

\clearpage
\begin{table}[!t]
\centering
\caption{Accounting of Inferential Claims}
\label{tab:inferential_accounting}
\renewcommand{\arraystretch}{1.4}
\setlength{\tabcolsep}{8pt}
\small
\begin{tabular}{@{}>{\raggedright\arraybackslash}p{3.0cm} >{\raggedright\arraybackslash}p{6.8cm} >{\raggedright\arraybackslash}p{4.8cm}@{}}
\toprule
\textbf{Evidentiary Status} & \textbf{Claim} & \textbf{Key Caveat} \\
\midrule
\multirow{3}{=}{Demonstrated} 
 & Moderate cross-state borrowing rather than wholesale copying or independent drafting; coalition asymmetry: protective $\beta = -0.076$ ($p = 0.022$); restrictive $\beta = -0.017$ ($p = 0.482$) & --- \\
 \addlinespace[6pt]
 & Organizational infrastructure and coalition membership predict similarity; contiguity null & Predicts pattern, not channel of transmission \\
 \addlinespace[6pt]
 & Zero common PAC donors across highest-similarity pairs; documentary corroboration of organizational co-presence & Confirms presence at case level, not drafting process \\
\midrule
\multirow{3}{=}{Established by inference} 
 & Reactive template mobilization as diffusion mechanism & Three rival mechanisms uneliminated (\S6.2.1) \\
 \addlinespace[6pt]
 & Vertical versus horizontal coalition architecture & Inferred from coefficient behavior \\
 \addlinespace[6pt]
 & Professionalism as capacity for adoption rather than innovation & Requires drafting-process evidence \\
\midrule
\multirow{2}{=}{Outside inferential reach} 
 & Internal drafting processes within legislative offices and advocacy organizations & Requires qualitative inquiry (\S\ref{subsec:future-qualitative}) \\
 \addlinespace[6pt]
 & Policy quality, efficacy, or durability of enacted legislation & Requires evaluative criteria outside text-based design \\
\bottomrule
\end{tabular}
\end{table}
\clearpage

The preceding sections have itemized specific methodological, specification-level, and scope limitations. This section consolidates them into an explicit accounting of the claims the computational evidence supports, the claims it establishes by inference, and the claims that lie outside its inferential reach. Table 6.1 summarizes the accounting; the discussion that follows develops each row in turn.

\textbf{What the design demonstrates.} The computational and documentary evidence assembled in this thesis establishes five claims with high confidence. First, the 2023 post-\emph{Dobbs} abortion corpus exhibits a characteristic textual structure: moderate baseline novelty (mean = 0.769, SD = 0.122) with a heavy central concentration between 0.70 and 0.90, consistent with widespread adaptation of shared frameworks rather than either wholesale copying or independent drafting. Second, that structure is coalition-asymmetric: protective bills are significantly less novel than the neutral baseline ($\beta = -0.076$, $p = 0.022$), while restrictive bills are statistically indistinguishable from it ($\beta = -0.017$, $p = 0.482$). Third, cross-state textual similarity is predicted by shared organizational infrastructure (PAC Jaccard $\beta = +0.039$, $p < 0.001$) and shared protective coalition membership ($\beta = +0.037$, $p < 0.001$), not by geographic contiguity ($\beta = +0.003$, $p = 0.734$). Fourth, the highest-similarity cross-state pairs exhibit near-zero direct financial connections between the legislators who sponsored the converging bills, establishing the decoupling of financial transfer from textual coordination at the case level. Fifth, documentary evidence from LegiScan legislative records, FEC contribution data, and advocacy organization publications corroborates the organizational interpretation by confirming the concurrent advocacy presence of the organizations whose Jaccard overlap predicts convergence.

\textbf{What the design establishes by inference.} Three claims are supported by the pattern evidence and the documentary record but rest on inferential extension rather than direct demonstration. First, the reactive template mobilization mechanism---the claim that pre-positioned templates were activated through affiliate coordination in response to the \emph{Dobbs} shock---is the most parsimonious account consistent with the observed signature but is not directly observed. The three rival mechanisms enumerated in \S\ref{subsec:text-necessary} remain uneliminated. Second, the vertical--horizontal architecture distinction between restrictive and protective coalitions is inferred from the differential behavior of the Both Restrictive and Both Protective coefficients across dyadic model specifications, not from direct observation of organizational drafting processes. The inference is theoretically grounded and empirically supported, but the organizational characterization---star-shaped versus mesh-shaped---is a theoretical construct applied to statistical patterns rather than a directly measured structural property. Third, the capacity-for-adoption reinterpretation of the professionalism finding is the most plausible reading of the negative $\beta$ in the bill-level regression, but the reinterpretation would be substantially strengthened by direct evidence about how legislative staff in high- and low-professionalism states access or fail to access externally produced templates.

\textbf{What the design cannot demonstrate.} Two categories of claims lie outside the design's inferential reach. First, mechanism-level claims about the internal drafting processes of legislative offices and advocacy organizations---who drafted what, which templates circulated through which channels, at what point specific phrasings entered the statutory text---require direct evidence the computational design cannot supply. The thesis establishes that the textual patterns are consistent with the organizational mechanisms Chapter~\ref{ch:discussion} theorizes; it does not establish that those mechanisms are the ones that actually produced the patterns. Second, evaluative claims about the relative quality, efficacy, or durability of the legislation the corpus contains are outside the scope of the text-based measure. Novelty scores describe textual distance, not policy merit.

\section{Future Research}
\label{sec:future-research}

The limitations documented in the preceding sections collectively define a research agenda organized around five extensions, listed in approximate order of priority for testing the framework's mechanism-level claims. Qualitative investigation of drafting processes (Section~\ref{subsec:future-qualitative}) is the most direct test of the reactive template mobilization mechanism and is positioned accordingly. Archival enrichment of the organizational measures (§6.6.2) and longitudinal extension to subsequent legislative cycles (§6.6.3) address measurement granularity and temporal scope. Cross-domain application (§6.6.4) and semantic methodological refinement (§6.6.5) test the framework's generalizability and the robustness of the textual measures to alternative representational strategies.

The extensions are listed in approximate order of priority for testing the framework's mechanism-level claims: qualitative investigation of drafting processes (\ref{subsec:future-qualitative}) is the most direct, archival enrichment and longitudinal extension (\ref{subsec:future-archival} - \ref{subsec:future-longitudinal}) address measurement and temporal scope, and cross-domain and methodological extensions (\ref{subsec:future-cross-domain} - \ref{subsec:future-methodological}) test generalizability and representational robustness.

\subsection{Qualitative Investigation of Legislative Drafting Processes}
\label{subsec:future-qualitative}

The most consequential extension would address the principal inferential gap identified throughout this chapter: the distance between the textual coordination pattern and the organizational mechanism theorized to produce it. Closing that gap requires direct evidence about the drafting processes internal to legislative offices and advocacy organizations---evidence that the computational and documentary materials assembled for this thesis cannot provide.

A natural next-stage research design would conduct semi-structured interviews with legislative staff, committee counsel, and advocacy organization personnel involved in drafting 2023 abortion legislation. The research questions such interviews would address are specific: Did staff in State A consult model statutory language circulated by Organization X? Did Organization X circulate such language, and if so, through what channels---direct staff contact, affiliate working groups, publicly accessible model-bill libraries, or informal networks? At what point in the drafting cycle did the converging phrasings identified by the computational analysis enter the statutory text? Were the converging phrasings present in the earliest available drafts, or did they emerge through amendment during the legislative process?

These questions are answerable. The personnel who drafted the 2023 legislation are identifiable through legislative records, and the organizational staff who supported the drafting process are identifiable through lobby-disclosure filings and organizational publications. The interviews would not require access to privileged or confidential materials; they would require access to the recollections and professional judgments of participants in a public legislative process. The methodological infrastructure for such inquiry is well established in the qualitative state-politics literature \parencite{hertel-fernandezStateCapture2019}.

This extension represents the natural continuation of the research program initiated here and is positioned as a priority for subsequent investigation. The computational analysis provides the map of convergence; the qualitative inquiry would supply the mechanism that produced it.

\subsection{Archival Enrichment of Organizational Measures}
\label{subsec:future-archival}

The binary Jaccard measure of organizational co-presence, while adequate for establishing the main effect documented in the dyadic regression, discards information about the depth and character of organizational embedding that richer data sources could supply. Three archival extensions would substantially improve the operationalization.

First, state-level lobby-disclosure filings would identify which organizations maintained registered lobbyists in each state during the 2023 legislative session, distinguishing active legislative engagement from passive electoral presence. Second, IRS Form 990 data on organizational staffing, program expenditures, and affiliated entities would provide a continuous measure of organizational depth that the binary Jaccard construction discards. Third, legislative liaison records---where states maintain them---would identify the specific organizational contacts through which model legislation and drafting assistance were channeled to legislative offices. Combining these sources with the DIME-derived Jaccard index would produce a multidimensional measure of organizational infrastructure whose components could be entered separately in the dyadic specification, testing which dimensions of organizational presence---electoral, lobbying, drafting, or clinical---predict the convergence signature.

The data-collection costs of this extension are nontrivial. Lobby-disclosure standards vary substantially across states, and some states' records are not digitized or standardized in ways that permit systematic cross-state comparison. IRS data are publicly available but require substantial cleaning and entity-matching work to link national organizations to their state-level affiliates. These costs are manageable within a funded research project but exceed the scope of the current thesis.

\subsection{Longitudinal Extension and Enactment-Timing Models}
\label{subsec:future-longitudinal}

Extending the corpus to include 2024 and 2025 enacted legislation would address two of the most consequential limitations identified in this chapter. First, it would test the constitutional vacuum diffusion framework's central temporal prediction: that the coalition asymmetry should attenuate as protective coalitions move beyond initial codification and restrictive coalitions develop new litigation-tested templates for the expanded policy frontier. A finding that protective novelty scores rose between 2023 and 2025 while restrictive scores fell would constitute strong confirmatory evidence for the framework; a finding that the asymmetry persisted or deepened would require theoretical revision.

Second, a multi-year corpus would supply the enactment-timing data the current design lacks. By incorporating bill introduction dates, committee hearing dates, floor vote dates, and gubernatorial signature dates across multiple legislative cycles, a longitudinal extension could construct event-history models that test directional diffusion hypotheses directly. The Vermont hub question---whether Vermont SB~37 functioned as a template source or merely as the textual center of a distributed co-drafting process---is answerable with introduction-date data that the current cross-sectional design does not exploit. More broadly, enactment-timing models would permit the exporter--importer--reinforcer typology to be grounded in observed temporal sequence rather than inferred from similarity patterns alone.

The longitudinal extension would also substantially improve statistical power for the interaction tests that the current dyadic sample cannot support (\S\ref{subsec:power-limits}). A three-year corpus, even under conservative assumptions, would roughly triple the dyadic sample and enable conditional analyses---whether the Both Protective effect varies with organizational density, whether professionalism moderates coalition effects---that the current 1{,}035-dyad sample renders underpowered.

\subsection{Cross-Domain Application}
\label{subsec:future-cross-domain}

The constitutional vacuum diffusion framework specifies scope conditions---predictable preemption collapse, high salience, dense bilateral advocacy infrastructure, morality-policy characteristics---under which its predictions should hold. Several candidate domains meet enough of these conditions to warrant systematic application, though each introduces complications the abortion case does not.

Firearms regulation following \emph{New York State Rifle \& Pistol Association v.\ Bruen} (2022) is the nearest analogue: a Supreme Court decision that altered the constitutional treatment of a highly salient, heavily organized policy domain with dense advocacy infrastructure on both sides. The key difference is that \emph{Bruen} narrowed permissible regulation rather than vacating a half-century of doctrinal infrastructure, producing an asymmetric shock that expanded the restrictive coalition's frontier without simultaneously compelling the protective coalition to codify previously assumed protections. The framework's predictions about bilateral mobilization would require modification for this asymmetric case.

Voting rights following \emph{Shelby County v.\ Holder} (2013) offers a partially completed natural experiment: the preclearance regime's partial invalidation opened a policy frontier for restrictive voting legislation, and the resulting legislative activity has been documented descriptively but not, to date, analyzed through text-as-data diffusion methods. The morality-policy characteristics are weaker---voting regulation involves substantial technical implementation alongside value conflicts---suggesting that hybrid diffusion mechanisms combining template adoption with local customization may predominate.

A hypothetical reversal of \emph{Obergefell v.\ Hodges} would closely parallel \emph{Dobbs} on most structural dimensions---predictable shock, morality policy, prepared organizational infrastructure, bilateral opportunities---but the organizational landscape differs in ways that could produce different coalition architectures. Whether these differences are consequential for the framework's predictions is an empirical question that pre-positioning the analytical infrastructure now would enable researchers to answer rapidly if the constitutional conditions materialize.

Each of these domains would test the framework's generalizability beyond the case that generated it. The framework's value for the broader diffusion literature depends on whether the structural conditions it identifies---rather than the substantive features of abortion politics---are what produce the observed diffusion dynamics.

\subsection{Methodological Extension: Semantic Embeddings}
\label{subsec:future-methodological}

The TF--IDF representation adopted for this thesis captures surface linguistic similarity but discards semantic equivalence across different wordings (\S\ref{subsec:tfidf-limits}). Replication of the core analyses using transformer-based embeddings---Legal-BERT \parencite{chalkidis-etal-2020-legal,devlinBERTPretrainingDeep2019} or domain-adapted successors---would test the robustness of the findings to a more semantically sensitive representation. The specific questions such replication would address are whether the coalition asymmetry, the Vermont hub pattern, and the Nevada--Washington near-replication persist when measured through embeddings that register functional equivalence rather than lexical overlap; and whether additional convergence invisible to n-gram methods emerges, particularly among bills that achieve similar doctrinal effects through different statutory language.

The expectation, given the strength of the surface-level convergence the TF--IDF analysis already detects, is that the substantive findings would remain stable under semantic replication. If that expectation is confirmed, the result would strengthen confidence that the patterns documented in Chapter~\ref{ch:findings} reflect genuine substantive coordination rather than artifacts of the representational method. If the expectation is disconfirmed---if, for instance, the coalition asymmetry disappears under semantic embeddings because restrictive bills achieve functional convergence through varied surface language---the finding would require substantial theoretical revision.

\section{Designed Scope, Not Incomplete Evidence}
\label{sec:designed-scope}

The limitations documented in this chapter are real, and the future research agenda they define is substantive. Neither observation diminishes the contributions the current design makes. The computational analysis provides the most systematic description to date of textual convergence patterns in post-\emph{Dobbs} state abortion legislation; the dyadic regression identifies organizational co-presence and coalition alignment as the structural predictors of that convergence while ruling out geographic proximity; and the documentary evidence corroborates the organizational interpretation at the case level. These are the questions the design was built to answer, and it answers them rigorously.

The questions it brackets---the internal organizational mechanisms through which specific statutory language was developed and circulated, the temporal sequence of template deployment within the 2023 cycle, the durability of the coalition asymmetry over subsequent legislative years, and the framework's generalizability to other policy domains---are questions that require different instruments. Qualitative inquiry into drafting processes, archival enrichment of organizational measures, longitudinal corpus extension, cross-domain replication, and semantic methodological refinement each address specific gaps the current design leaves open, and each is feasible within the infrastructure the thesis has established.

The distinction between designed scope and incomplete evidence matters for how the thesis's contributions should be assessed. A research design that attempts to answer every question answers none of them well. The choice to prioritize computational text analysis over qualitative mechanism-tracing, to analyze a single legislative year rather than a multi-year panel, and to focus on one policy domain rather than several was a methodological prioritization that enabled the precision the findings achieve. The limitations this chapter documents are the costs of that prioritization, and they are specified here not as concessions but as the boundary conditions within which the findings hold and beyond which subsequent investigation begins.

What remains, then, is to consolidate what the thesis has established within those boundaries. The preceding chapters have assembled the empirical record, interpreted its theoretical implications, and specified the limits of the interpretive claims. Chapter~7 draws together these strands: it restates the central contributions in their most concise form, locates the thesis within the broader trajectory of diffusion scholarship, and closes with the question that the 2023 record poses most urgently for the field---whether the patterns I document represent a transitional anomaly or the first evidence of a structural transformation in how American states make policy after constitutional shocks.